\chapter*{Proto Class}
\phantomsection
\addstarredchapter{Proto Class}
\fancyhead[LE]{\thepage} 
\fancyhead[RE]{Proto Class}
\fancyhead[RO]{\thepage} 
\fancyhead[LO]{Proto Class}
\fancyfoot[C]{\thepage}

\renewcommand{\headrulewidth}{0.1pt}

\setlength{\parindent}{0cm}
\setlength{\parskip}{9pt}
\vspace{-1cm}

% -------------------------------------------------------------
Convention d'ecriture Capital = Folders, camelCase filenames. 
Structuration des données d'entrée pour \hydrot\ :

\begin{verbatim}
idCard.md
temporalSupport.md
COMPARTMENTi/
....

\end{verbatim}


Contenu de \texttt{COMPARTMENTi} :
\begin{verbatim}
idCard.md
SPATIAL\_SUPPORT (GeoData)
PARAMETERS (GeoData)
FORCINGS (NONE if internal compartment)
DATA (OBS)
SIMUL (Simulation outputs)
....

\end{verbatim}

Où \texttt{idCard.md}
\begin{verbatim}
   modelName = 
   version = 
   repo-path = 
   branch = (if void use commit_hash)
   commit_hash = (if void, use branch)
   OS = (name,version)
   compilo = (gcc,version),(gfortran,version), autres ....
   #scientific_certificate = none/dirty/ok %calculé à l'initialisation de HydroTwin
   epsilon = (distance maximum absolue entre hypercubes)

\end{verbatim}


Chaque FOLDER contient l'information nécessaire pour init \texttt{SimplifiedGitSynchroRecord} (section \ref{p2c2_ssGitRec}, p. \pageref{p2c2_ssGitRec})


Si \texttt{MODELID} existe il initialise \texttt{COMPARTMENTi/MODELID}